% Unit 105 terjemahan Bahasa Indonesia dari chapter7.tex baris 2945--3002.
% Sumber: Methods of Algebra, Volume 2, commit
% 9a5803ff77dd3257484cb177f851a73770a59dd3 (CC BY 4.0).
% stable-unit-id: o014.aljabr2.chapter7.exercises-c
% Cakupan: latihan penutup Bab 7 tentang penurunan Galois, pemuntiran
% kohomologi nonabelian, dan bentuk kuadratik; menutup lingkungan Exercises.

\hypertarget{o014.aljabr2.chapter7.exercises-c}{ }
% segment-id: o014.aljabr2.chapter7.exercises-c.ex001
	\item Mengikuti penjelasan pada akhir \S\ref{sec:descent-Galois}, berikan
	bukti langsung Teorema Penurunan Galois~\ref{prop:Galois-descent} melalui
	langkah-langkah berikut. Gunakan notasi dari bagian tersebut: misalkan
	$M\in\Obj(\cate{Vect}^{\Gamma,\infty}(L))$, dan definisikan pemetaan
	kanonik $\iota_M$ melalui \eqref{eqn:Galois-descent-iotaM}.
	\begin{enumerate}[(i)]
% segment-id: o014.aljabr2.chapter7.exercises-c.i001
		\item Misalkan $M'$ adalah subruang $\Gamma$-invarian dari
		$L\dotimes{K}(M^\Gamma)$, yakni $\Gamma M'\subset M'$. Buktikan bahwa
		jika $M'\cap(1\otimes M^\Gamma)=\{0\}$, maka $M'=\{0\}$.

% segment-id: o014.aljabr2.chapter7.exercises-c.hi001
		\begin{hint}
			Pilih suatu basis $(y_i)_{i\in I}$ dari $M^\Gamma$. Jika
			$M'\neq\{0\}$, pilih suatu ekspresi taknol di dalamnya dengan jumlah
			suku sesedikit mungkin, $\sum_{i\in I}\ell_i\otimes y_i$ (dengan
			jumlah hingga). Setelah normalisasi yang sesuai, kita dapat
			mengandaikan bahwa ekspresi itu
			berbentuk
% segment-id: o014.aljabr2.chapter7.exercises-c.q001
			\[ m' = 1 \otimes y_{i_0} + \ell_1 \otimes y_{i_1} + \cdots, \quad \ell_1 \notin K. \]
			Ambil $\sigma\in\Gamma$ sehingga
			$\sigma(\ell_1)\neq\ell_1$, lalu tinjau $m'-\sigma(m')$ untuk
			memperoleh kontradiksi.
		\end{hint}

% segment-id: o014.aljabr2.chapter7.exercises-c.i002
		\item Buktikan bahwa
		$\iota_M:L\dotimes{K}(M^\Gamma)\to M$ injektif.
% segment-id: o014.aljabr2.chapter7.exercises-c.hi002
		\begin{hint}
			Tinjau $M':=\Ker(\iota_M)$.
		\end{hint}

% segment-id: o014.aljabr2.chapter7.exercises-c.i003
		\item Buktikan bahwa $\iota_M$ surjektif.

% segment-id: o014.aljabr2.chapter7.exercises-c.hi003
		\begin{hint}
			Tinjau suatu subperluasan Galois hingga $E|K$ dari $L|K$. Ambil
			$a_1,\ldots,a_n$ sebagai basis $E$ selaku ruang vektor atas $K$, dan pilih
			wakil $\identity=\sigma_1,\ldots,\sigma_n\in\Gamma$ bagi unsur-unsur
			$\Gal(E|K)$. Dari teori medan diketahui bahwa
			$(\sigma_i(a_j))_{1\leq i,j\leq n}$ merupakan matriks invertibel di
			atas $E$; tuliskan inversnya sebagai
			$(b_{ij})_{1\leq i,j\leq n}$. Untuk
			$m\in M^{\Gal(L|E)}$, turunkan
% segment-id: o014.aljabr2.chapter7.exercises-c.q002
			\[ m = \sum_{i=1}^n b_{i1} \sum_{j=1}^n \sigma_j(a_i m) \; \in \Image(\iota_M). \]
		\end{hint}

% segment-id: o014.aljabr2.chapter7.exercises-c.i004
		\item Gunakan hasil-hasil di atas untuk menunjukkan bahwa
		\eqref{eqn:Galois-descent-adjunction} merupakan adjungsi yang sekaligus
		merupakan ekuivalensi,
		sehingga diperoleh bukti langsung
		Teorema~\ref{prop:Galois-descent}.
	\end{enumerate}

% segment-id: o014.aljabr2.chapter7.exercises-c.ex002
	\item Misalkan $G$ grup dan $L|K$ perluasan Galois medan. Tuliskan kategori
	modul-$G$ dengan koefisien dalam medan $K$ (atau $L$) sebagai
	$G\dcate{Mod}$ (atau $G\dcate{Mod}_L$); lihat
	Definisi~\ref{def:G-mod}. Gunakan metode penurunan Galois dari
	\S\S\ref{sec:descent-Galois}---\ref{sec:Galois-H1} untuk mendeskripsikan
	$G\dcate{Mod}$ melalui objek-objek $G\dcate{Mod}_L$ yang dilengkapi aksi
	$\Gal(L|K)$ yang sesuai.

% segment-id: o014.aljabr2.chapter7.exercises-c.ex003
	\item Gunakan notasi \S\ref{sec:Galois-H1}. Misalkan
	$c\in Z^1(\Gamma,G(\mathbf{t}_{0,L}))$. Dari bukti
	Teorema~\ref{prop:Galois-H1}, ingat bahwa melalui penurunan Galois, $c$
	mendefinisikan ruang vektor-$K$ $W$ beserta data $\mathbf{t}$ di atasnya,
	dan juga $h:\mathbf{t}_{0,L}\rightiso\mathbf{t}_L$.
	\begin{enumerate}[(i)]
% segment-id: o014.aljabr2.chapter7.exercises-c.i005
		\item Tunjukkan bahwa $g\mapsto h^{-1}gh$ memberikan isomorfisme grup
		$\nu:G(\mathbf{t}_L)\rightiso G(\mathbf{t}_{0,L})$. Tunjukkan bahwa
% segment-id: o014.aljabr2.chapter7.exercises-c.q003
		$ f \xmapsto{\sigma \in \Gamma} c(\sigma)({}^\sigma f)c(\sigma)^{-1} $
		memberikan aksi $\Gamma$ baru pada $G(\mathbf{t}_{0,L})$. Tuliskan
		struktur ini sebagai ${}^cG(\mathbf{t}_{0,L})$, dan tunjukkan bahwa
		$\nu:G(\mathbf{t}_L)\rightiso{}^cG(\mathbf{t}_{0,L})$ bersifat
		$\Gamma$-ekuivarian.

% segment-id: o014.aljabr2.chapter7.exercises-c.i006
		\item Tunjukkan bahwa rumus berikut memberikan pemetaan yang
		terdefinisi dengan baik:
% segment-id: o014.aljabr2.chapter7.exercises-c.q004
		\begin{align*}
			Z^1(\Gamma, G(\mathbf{t}_L)) & \to Z^1(\Gamma, G(\mathbf{t}_{0, L})) \\
			\tilde{c} & \mapsto \left[\sigma \mapsto \nu(\tilde{c}(\sigma)) c(\sigma) \right].
		\end{align*}
		Pemetaan ini menginduksi bijeksi
		$\Hm^1(\Gamma,G(\mathbf{t}_L))\to
		\Hm^1(\Gamma,G(\mathbf{t}_{0,L}))$ yang memetakan titik pangkal ke
		$[c]$.
% segment-id: o014.aljabr2.chapter7.exercises-c.hi004
		\begin{hint}
			Bandingkan dengan konstruksi ``pemuntiran'' dalam latihan
			CHref{sec:group-coh}.
		\end{hint}

% segment-id: o014.aljabr2.chapter7.exercises-c.i007
		\item Tunjukkan bahwa, di bawah korespondensi dalam
		Teorema~\ref{prop:Galois-H1}, bijeksi ini bersesuaian dengan kesamaan
		$\mathscr{T}(\mathbf{t})=\mathscr{T}(\mathbf{t}_0)$.
	\end{enumerate}

% segment-id: o014.aljabr2.chapter7.exercises-c.ex004
	\item Misalkan $K$ medan dengan $\mathrm{char}(K)\neq2$. Suatu bentuk
	bilinear simetris takdegenerat $q:V\times V\to K$ pada ruang vektor-$K$
	berdimensi hingga $V$ selanjutnya disebut bentuk kuadratik takdegenerat.
	Dalam grup ortogonal $\mathrm{O}(q)=\mathrm{O}(V,q)$, tuliskan subgrup
	yang ditentukan oleh $\det=1$ sebagai
	$\SO(q)=\SO(V,q)$. Untuk setiap perluasan Galois $L|K$, terdapat bentuk
	kuadratik takdegenerat $(V_L,q_L)$ beserta grup-grup terkait;
	semuanya dilengkapi aksi mulus dari $\Gamma:=\Gal(L|K)$.
	\begin{enumerate}[(i)]
% segment-id: o014.aljabr2.chapter7.exercises-c.i008
		\item Tuliskan $\mu_2:=\{\pm1\}\subset K^\times$. Tunjukkan bahwa
		$\mathrm{O}(q)/\SO(q)\simeq\mu_2$. Untuk hasil bagi
		$\mathrm{O}(q_L)/\SO(q_L)$, jelaskan hubungan isomorfisme ini
		dan aksi $\Gamma$.

% segment-id: o014.aljabr2.chapter7.exercises-c.i009
		\item Dengan menggunakan perangkat \S\ref{sec:Galois-H1}, tunjukkan
		bahwa unsur-unsur $\Hm^1(\Gamma,\mathrm{O}(q_L))$ berkorespondensi
		bijektif dengan kelas isomorfisme bentuk kuadratik takdegenerat
		$(V',q')$ di atas $K$ yang memenuhi $\dim V'=\dim V$, dengan titik
		pangkal bersesuaian dengan $(V,q)$.

% segment-id: o014.aljabr2.chapter7.exercises-c.i010
		\item Terapkan Teorema~\ref{prop:nonabelian-long} untuk memperoleh
		barisan eksak himpunan bertitik
% segment-id: o014.aljabr2.chapter7.exercises-c.q005
		\[ \mathrm{O}(q) \xrightarrow{\det} \mu_2 \to \Hm^1(\Gamma, \SO(q_L)) \xrightarrow{\varphi} \Hm^1(\Gamma, \mathrm{O}(q_L)) \xrightarrow{\psi} \Hm^1(\Gamma, \mu_2). \]
		Cobalah tunjukkan bahwa $\varphi$ injektif.
% segment-id: o014.aljabr2.chapter7.exercises-c.hi005
		\begin{hint}
			Surjektivitas $\det$ hanya menunjukkan bahwa
			$\varphi^{-1}(\text{titik pangkal})=\{\text{titik pangkal}\}$.
			Untuk menunjukkan injektivitas, kita harus memvariasikan $(V,q)$ dan
			menerapkan teknik pemuntiran dari latihan sebelumnya.
		\end{hint}

% segment-id: o014.aljabr2.chapter7.exercises-c.i011
		\item Tunjukkan bahwa
		$\Hm^1(\Gamma,\mu_2)\simeq K^\times/K^{\times,2}$. Tunjukkan pula bahwa
		$\psi$ sama dengan pemetaan
% segment-id: o014.aljabr2.chapter7.exercises-c.q006
		\[ \left[ (V', q'):\; \text{bentuk kuadratik takdegenerat} \right] \mapsto \mathrm{disc}(q') / \mathrm{disc}(q), \]
		Di sini, pilih sebarang basis $V'$, identifikasi $q'$ dengan
		suatu matriks simetris, lalu definisikan
% segment-id: o014.aljabr2.chapter7.exercises-c.q007
		\[ \mathrm{disc}(q') := (-1)^n \det(q') \bmod K^{\times 2}, \quad \dim V = 2n \;\text{atau}\; 2n+1. \]

% segment-id: o014.aljabr2.chapter7.exercises-c.i012
		\item Tunjukkan bahwa unsur-unsur
		$\Hm^1(\Gamma,\SO(q_L))$ berkorespondensi bijektif dengan kelas
		isomorfisme bentuk kuadratik takdegenerat $(V',q')$ di atas $K$ yang
		memenuhi $\dim V'=\dim V$ dan
		$\mathrm{disc}(q')=\mathrm{disc}(q)$, dengan titik pangkal
		bersesuaian dengan $(V,q)$.
	\end{enumerate}
% segment-id: o014.aljabr2.chapter7.exercises-c.c001
\end{Exercises}
